\subsection{李群的切向量群}

命题 4. 设 G 是李群。则 T(G) 配备与 G 的乘法相切的复合律后是李群。T(G) 的单位元是向量 \(0_{e}\)。

这由命题 1 和命题 2 得出。

命题 5. 设 G、H 是李群，f 是从 G 到 H 的态射。则 T(f) 是从李群 T(G) 到李群 T(H) 的态射。

我们知道 \(\mathrm{T}(f)\) 是解析的。另一方面，令 m（相应地，n）表示 G（相应地，H）上的乘法。则 \(f \circ m = n \circ (f \times f)\)，从而

\[
\mathrm{T} (f) \circ \mathrm{T} (m) = \mathrm{T} (n) \circ (\mathrm{T} (f) \times \mathrm{T} (f)),
\]

这表达了 \(T(f)\) 是群同态这一事实。

推论。设 \(G_{1},\ldots,G_{n}\) 是李群。从流形 \(\mathrm{T}(\mathrm{G}_{1}\times\cdots\times\mathrm{G}_{n})\) 到流形 \(\mathrm{T}(\mathrm{G}_{1})\times\cdots\times\mathrm{T}(\mathrm{G}_{n})\) 的典范同构是李群同构。

\(pr_{i}\) 是从 \(G_{1} \times \cdots \times G_{n}\) 到 \(G_{i}\) 的态射，因而 \(T(pr_{i})\) 是从 \(T(G_{1} \times \cdots \times G_{n})\) 到 \(T(G_{i})\) 的态射。

命题 6. 设 G 是李群。

\begin{enumerate}
\def\labelenumi{(\roman{enumi})}
\item
  典范投影 \(p: \mathrm{T}(\mathrm{G}) \to \mathrm{G}\) 是李群态射。
\item
  \(p\) 的核是 \(\mathrm{T}_e(\mathrm{G})\)。它是 \(\mathrm{T}(\mathrm{G})\) 的李子群。由 \(\mathrm{T}_e(\mathrm{G})\) 上由 \(\mathrm{T}(\mathrm{G})\) 上的结构诱导的李群结构，是完备赋范空间 \(\mathrm{T}_e(\mathrm{G})\) 的李群结构。
\item
  零截面 \(s\) 是李群 \(\mathbf{G}\) 到李子群 \(s(\mathbf{G})\) 的同构，该子群属于 \(\mathrm{T}(\mathbf{G})\)（我们将此子群识别为 \(\mathbf{G}\)）。
\item
  李群 T(G) 是 G 与 \(T_{e}(G)\) 的半直积。
\end{enumerate}

断言 (i) 由 (5) 得出。考虑命题 2 (ii)，断言 (ii) 是显然的。断言 (iii) 和 (iv) 由 (6) 及 § 1 的命题 8 得出。

令 \(u \in \mathrm{T}(\mathrm{G})\) 且 \(g \in \mathrm{G}\)。根据 (3) 和 (4)，乘积 \(ug\)、\(gu\)（在群 \(\mathrm{T}(\mathrm{G})\) 中计算）分别是 \(u\) 在 \(\mathrm{T}(\delta(g^{-1}))\) 和 \(\mathrm{T}(\gamma(g))\) 下的像。由 § 1 命题 17 的推论 2 可知，映射 \((g, u) \mapsto gu\) 从 \(\mathrm{G} \times \mathrm{T}_e(\mathrm{G})\) 到 \(\mathrm{T}(\mathrm{G})\) 是平凡向量丛 \(\mathrm{G} \times \mathrm{T}_e(\mathrm{G})\)（基空间为 \(G\)）到向量丛 \(\mathrm{T}(\mathrm{G})\) 的同构。其逆同构称为 \(\mathrm{T}(\mathrm{G})\) 的左平凡化。考察映射 \((g, u) \mapsto ug\)，即可类似地定义 \(\mathrm{T}(\mathrm{G})\) 的右平凡化。

命题 7. 设 G 是李群，M 是类为 \(\mathbf{C}^r\) 的流形，\(f\) 和 \(g\) 是从 M 到 G 的类为 \(\mathbf{C}^r\) 的映射，从而 \(fg\) 是从 M 到 G 的类为 \(\mathbf{C}^r\) 的映射。令 \(m \in \mathbf{M}\)，\(x = f(m)\)，\(y = g(m)\)，\(u \in \mathrm{T}_m(\mathbf{M})\)。则

\[
(\mathrm{T} f g) u = \mathrm{T} (f) u. y + x. \mathrm{T} (g) u.
\]

令 \(m\) 表示 G 的乘法。则 \(fg = m \circ (f, g)\)。现在

\[
\mathrm{T} (f, g) (u) = (\mathrm{T} (f) u, \mathrm{T} (g) u),
\]

因此 \(\mathrm{T}(fg)u = \mathrm{T}(f)u.\mathrm{T}(g)u\)。于是只需在 (1) 中将 \(f\) 替换为 \(m\)。

推论。令 \(n \in \mathbf{Z}\)。在 \(e\) 处，映射 \(g \mapsto g^n\)（从 \(G\) 到 \(G\)）的切映射是映射 \(x \mapsto nx\)，从 \(T_e(G)\) 到 \(T_e(G)\)。

对于 \(n \geqslant 0\)，这由命题 7 对 \(n\) 作归纳得到。另一方面，在 \(e\) 处到映射 \(g \mapsto g^{-1}\) 的切映射是映射 \(x \mapsto -x\)（第 1 号，命题 2）。

设 G 是李群，X 是类为 \(C^r\) 的流形，且 \((g, x) \mapsto gx\) 是 G 在 X 上的类为 \(C^r\) 的左作用律。仿照命题 1，可得到 T(G) 在 T(X) 上的类为 \(C^{r-1}\) 的左作用律，我们仍将其记为 \((u, v) \mapsto uv\)。将 G（相应地，X）识别为 T(G)（相应地，T(X)）的零截面的像，由 (6) 可见，T(G) 在 T(X) 上的左作用律延拓了 G 在 X 上的左作用律。对于所有 \(u \in \mathrm{T}_{g}(G)\) 和 \(v \in \mathrm{T}_{x}(X)\)，根据 (1)，

\[
u v = g v + u x.\tag{7}
\]

若 \(g \in G\) 且 \(v \in T_x(X)\)，则根据 (3)，\(gv\) 是 \(v\) 在 \(x\) 处对映射 \(y \mapsto gy\) 的切映射下的像，该映射从 \(X\) 到 \(X\)。这个切映射是从 \(T_x(X)\) 到 \(T_{gx}(X)\) 的同构。特别地，

\[
g (v + v ^ {\prime}) = g v + g v ^ {\prime}, \quad g (\lambda v) = \lambda (g v) \quad \text { for } v, v ^ {\prime} \text { in } T _ {x} (X), \lambda \in K.\tag{8}
\]

若 \(x \in \mathbf{X}\) 且 \(u \in \mathrm{T}_g(\mathrm{G})\)，则根据 (4)，\(ux\) 是 \(u\) 在 \(g\) 处对映射 \(h \mapsto hx\) 的切映射下的像，该映射从 \(\mathbf{G}\) 到 \(\mathbf{X}\)。因此

\[
(9) \quad (u + u ^ {\prime}) x = u x + u ^ {\prime} x, \qquad (\lambda u) x = \lambda (u x) \quad \text { for } u, u ^ {\prime} \text { in } \mathrm{T} _ {g} (\mathrm{G}), \lambda \in \mathrm{K}.
\]

上述结果可应用于李群通过左（相应地，右）平移作用于自身的情形。T(G) 在 T(G) 上相应的作用律由李群 T(G) 的左（相应地，右）平移定义。因此公式 (7)、(8) 和 (9) 在 T(G) 中成立。

命题 8. 设 \(\mathbf{G}_1\) 和 \(\mathbf{G}_2\) 是李群，\(\mathbf{X}_1\) 和 \(\mathbf{X}_2\) 是类为 \(\mathbf{C}^r\) 的流形，\(f_i\) 是类为 \(\mathbf{C}^r\) 的左作用律，作用于 \(\mathbf{G}_i\) 和 \(\mathbf{X}_i\)（\(i = 1, 2\)）。令 \(\phi\) 为从 \(\mathbf{G}_1\) 到 \(\mathbf{G}_2\) 的态射，\(\psi\) 为 \(\phi\)-态射，它从 \(\mathbf{X}_1\) 到 \(\mathbf{X}_2\)。则 \(\mathrm{T}(\psi)\) 是 \(\mathrm{T}(\phi)\)-态射，它从 \(\mathrm{T}(\mathbf{X}_1)\) 到 \(\mathrm{T}(\mathbf{X}_2)\)。

\(f_{2}\circ (\phi \times \psi) = \psi \circ f_{1},\) 从而

\[
\mathrm{T} (f _ {2}) \circ (\mathrm{T} (\phi) \times \mathrm{T} (\psi)) = \mathrm{T} (\psi) \circ \mathrm{T} (f _ {1}).
\]

设 G 是李群，X 是类为 \(C^{r}\) 的流形，且 \((g, x) \mapsto gx\) 是 G 在 X 上的类为 \(C^{r}\) 的左作用律。令 I 为包含 0 的 K 的开子集，\(\gamma: I \to G\) 为类为 \(C^{r}\) 的映射，满足 \(\gamma(0) = e\)。令

\[
a = \mathrm{T} _ {0} (\gamma) 1 \in \mathrm{T} _ {e} (\mathrm{G}).
\]

令 \(x \in X\)。根据 (4)，\(ax\) 是 I 在 0 处的切向量 1 在映射 \(\lambda \mapsto \gamma(\lambda)x\) 的切映射下的像。因此，向量场 \(x \mapsto ax\) 定义在 \(X\) 上，是由映射 \((\lambda, x) \mapsto \gamma(\lambda)x\) 定义的向量场，其含义见《Differentiable and Analytic Manifolds》R, 8.4.5。
% semantic-fragments: legacy-input-seam
\relax{}
